% Vocabulary \& Definitions

\begin{description}[leftmargin=0.5cm, style=nextline, noitemsep, topsep=0pt]
    \item[Arbitrary] Bất kỳ, không ràng buộc.
    \item[Bipartite Graph] Đồ thị 2 phía (cạnh chỉ nối giữa 2 phía).
    \item[Matching] Cặp ghép (tập cạnh không chung đỉnh).
    \item[Maximal Matching] Cặp ghép tối đại (không thể thêm cạnh).
    \item[Maximum Matching] Cặp ghép cực đại (nhiều cạnh nhất).
    \item[Perfect Matching] Cặp ghép hoàn hảo (phủ kín mọi đỉnh).
    \item[Independent Set] Tập độc lập (tập đỉnh không có cạnh nối).
    \item[Vertex Cover] Phủ đỉnh (tập đỉnh chạm vào mọi cạnh).
    \item[Edge Cover] Phủ cạnh (tập cạnh chạm vào mọi đỉnh).
    \item[Lexicographically] Thứ tự từ điển (so sánh như chuỗi).
    \item[Monotonicity] Tính đơn điệu (luôn tăng/giảm $\rightarrow$ chặt nhị phân).
    \item[Idempotent] Lũy đẳng ($f(x, x) = x \rightarrow$ GCD, Max, Min).
    \item[Mutually exclusive] Xung khắc (không thể xảy ra cùng lúc).
    \item[Complement] Phần bù / Đồ thị bù (đảo ngược cạnh).
    \item[Inversion] Cặp nghịch thế ($i < j$ nhưng $A[i] > A[j]$).
    \item[Adjacent] Kề nhau (đỉnh kề, ô lưới kề cạnh/góc).
    \item[Consecutive] Liên tiếp (phần tử liên tiếp trong mảng).
    \item[Alternative] Luân phiên / Thay thế (ví dụ: chuỗi ký tự dạng 010101).
    \item[Parity] Tính chẵn lẻ (bit parity, tổng chẵn/lẻ).
    \item[Coprime] Nguyên tố cùng nhau ($\text{gcd}(a,b) = 1$).
    \item[Bijective] Song ánh (ánh xạ 1-1 tương ứng hoàn toàn).
    \item[Invariants] Đại lượng bất biến (giá trị không đổi qua các thao tác).
    \item[Predecessor / Successor] Phần tử đứng ngay trước / đứng ngay sau.

    \item[With replacement] Có hoàn lại (bốc thăm xong bỏ lại vào hộp $\rightarrow$ xác suất không đổi).
    \item[Without replacement] Không hoàn lại (bốc thăm xong giữ luôn bên ngoài $\rightarrow$ xác suất thay đổi).
    \item[Composite number] Hợp số (số có nhiều hơn 2 ước dương).
    \item[Factor / Divisor / Multiple] Ước số / Ước số / Bội số.
    \item[Quotient / Remainder] Thương số / Số dư trong phép chia.
    \item[Coefficient] Hệ số (ví dụ: hệ số của đa thức $x^k$).
    \item[Modulo / Congruence] Phép lấy số dư / Quan hệ đồng dư ($a \equiv b \pmod m$).

    \item[Collinear] Thẳng hàng (các điểm cùng nằm trên 1 đường thẳng).
    \item[Coplanar] Đồng phẳng (các điểm cùng thuộc 1 mặt phẳng).
    \item[Convex / Concave polygon] Đa giác lồi / Đa giác lõm.
    \item[Convex Hull] Bao lồi (đa giác lồi nhỏ nhất bao quanh tập điểm).
    \item[Cross product] Tích có hướng (tìm hướng quay CCW/CW, tính diện tích).
    \item[Dot product] Tích vô hướng (tính góc, độ dài hình chiếu).
    \item[Perpendicular / Orthogonal] Vuông góc / Trực giao ($\text{dot product} = 0$).
    \item[Counter-clockwise (CCW)] Ngược chiều kim đồng hồ ($\text{cross product} > 0$).
    \item[Clockwise (CW)] Thuận chiều kim đồng hồ ($\text{cross product} < 0$).
    \item[Degenerate polygon] Đa giác suy biến (các đỉnh thẳng hàng hoặc trùng nhau).
    \item[Bounding box] Khung hình chữ nhật nhỏ nhất bao quanh vật thể.
    \item[Incircle] Đường tròn nội tiếp (tiếp xúc trong với các cạnh của đa giác).
    \item[Incenter] Tâm đường tròn nội tiếp (giao điểm các đường phân giác trong).
    \item[Inradius ($r$)] Bán kính đường tròn nội tiếp (Tính nhanh: $r = \text{Area}/s$ với $s$ là nửa chu vi).
    \item[Circumcircle] Đường tròn ngoại tiếp (đi qua tất cả các đỉnh của đa giác).
    \item[Circumcenter] Tâm đường tròn ngoại tiếp (giao điểm các đường trung trực).
    \item[Circumradius ($R$)] Bán kính đường tròn ngoại tiếp (Tam giác: $R = abc / (4 \cdot \text{Area})$).
    \item[Excircle / Excenter] Đường tròn bàng tiếp / Tâm bàng tiếp (tiếp xúc 1 cạnh và 2 cạnh kéo dài).
    \item[Concyclic] Đồng viên (các điểm cùng nằm trên một đường tròn).
    \item[Tangent] Tiếp tuyến (đường thẳng tiếp xúc đường tròn tại 1 điểm, vuông góc bán kính).
    \item[Secant] Cát tuyến (đường thẳng cắt đường tròn tại 2 điểm phân biệt).
    \item[Chord] Dây cung (đoạn thẳng nối 2 điểm trên đường tròn).
    \item[Diameter / Radius] Đường kính / Bán kính.
    \item[Centroid] Trọng tâm (giao điểm các đường trung tuyến; hoặc trung bình cộng tọa độ đa giác).
    \item[Orthocenter] Trực tâm (giao điểm các đường cao của tam giác).
    \item[Sector / Segment] Hình quạt tròn (giới hạn bởi 2 bán kính và cung) / Hình viên phân (giới hạn bởi dây cung và cung).

    \item[Incidence Matrix] Ma trận liên thuộc (biểu diễn mối quan hệ giữa đỉnh và cạnh).
    \item[Identity Matrix] Ma trận đơn vị (đường chéo chính bằng 1, các vị trí khác bằng 0).
    \item[Transpose] Ma trận chuyển vị (đổi hàng thành cột).
    \item[Determinant] Định thức của ma trận (dùng tính diện tích đa giác, giải hệ PT, Matrix-Tree theorem).
    \item[Planar Graph] Đồ thị phẳng (đồ thị có thể vẽ trên mặt phẳng mà không có hai cạnh nào cắt nhau).
    \item[Self-loop / Multiple edges] Cạnh tự khuyên (nối đỉnh với chính nó) / Đa cạnh (nhiều cạnh nối cùng một cặp đỉnh).
    \item[Simultaneously] Đồng thời cùng một lúc (thao tác diễn ra đồng loạt).
    \item[Respectively] Lần lượt, theo thứ tự tương ứng (ví dụ: các mảng $A, B$ có độ dài $X, Y$ tương ứng).
    \item[Sufficient / Necessary] Điều kiện đủ / Điều kiện cần.
    \item[Trivial / Non-trivial] Hiển nhiên, tầm thường (ví dụ: trường hợp rỗng, bằng 0) / Không hiển nhiên.
    \item[Symmetric / Asymmetric] Đối xứng / Bất đối xứng.
    \item[Impartial / Partisan game] Trò chơi công bằng (mọi người chơi cùng tập nước đi) / Không công bằng (tập nước đi khác nhau).

    \item[Chordal Graph] Đồ thị dây cung (mọi chu trình độ dài $\ge 4$ đều có một cạnh nối giữa 2 đỉnh không kề nhau).
    \item[Tournament] Đồ thị giải đấu (đồ thị có hướng thu được bằng cách định hướng cho mọi cạnh của một đồ thị đầy đủ).
    \item[Self-complementary] Tự bù (đồ thị đẳng cấu với chính đồ thị bù của nó).
    \item[Eulerian / Hamiltonian] Đồ thị Euler (có chu trình qua mọi cạnh) / Đồ thị Hamilton (có chu trình qua mọi đỉnh).
    \item[Bridge-connected] Liên thông qua cầu (thành phần không có đỉnh khớp nhưng vẫn chứa cầu).
    \item[Topologically equivalent] Đồng phôi (hai cấu trúc đồ thị có thể biến đổi qua lại bằng cách thêm/bớt đỉnh bậc 2).

    \item[Submask / Supermask] Mặt nạ con (tập con bit) / Mặt nạ mẹ (tập cha bit) $\rightarrow$ Dấu hiệu nhận biết bài toán SOS DP.
    \item[Subadditivity / Superadditivity] Tính dưới cộng tính ($f(x+y) \le f(x) + f(y)$) / Tính trên cộng tính ($f(x+y) \ge f(x) + f(y)$).
    \item[Submodularity] Tính dưới mô-đun ($f(A \cup \{x\}) - f(A) \ge f(B \cup \{x\}) - f(B)$ với $A \subseteq B$ $\rightarrow$ thuật toán tham lam đạt tối ưu).
    \item[Quasiconvex / Quasiconcave] Hàm tựa lồi (tập mức dưới là tập lồi) / Hàm tựa lõm (tập mức trên là tập lồi) $\rightarrow$ Chặt tam phân.
    \item[Injective / Surjective] Đơn ánh (không có 2 phần tử trùng ảnh) / Toàn ánh (mọi phần tử đích đều có tạo ảnh).
    \item[Asymptotic complexity] Độ phức tạp tiệm cận (ước lượng Big-O khi dữ liệu tiến ra vô cùng).
    \item[Non-decreasing / Non-increasing] Không giảm ($\ge$, có thể bằng nhau) / Không tăng ($\le$, có thể bằng nhau) $\rightarrow$ Ngặt hơn strictly.

    \item[Collinear \& Non-coincident] Thẳng hàng nhưng không trùng nhau (bẫy hình học cực nặng khi xét giao điểm).
    \item[Strictly convex] Lồi ngặt (không có 3 điểm nào trên biên của đa giác lồi thẳng hàng).
    \item[Orthogonal polygon] Đa giác trực giao (tất cả các cạnh chỉ nằm ngang hoặc thẳng đứng, góc trong luôn là $90^\circ$ hoặc $270^\circ$).
    \item[Self-intersecting] Tự cắt (đa giác có các cạnh cắt nhau không phải tại đỉnh $\rightarrow$ không áp dụng được Shoelace thông thường).
    \item[Minkowski Sum] Tổng Minkowski (tập hợp tổng các vector của hai tập điểm $\rightarrow$ dùng để gộp nhanh hai bao lồi).

    \item[Ad-hoc] Bài toán không có thuật toán chuẩn mực (chỉ dựa vào quan sát, mẹo, hoán vị hoặc cấu trúc đặc biệt để giải).
    \item[Amortized time] Độ phức tạp thời gian khấu hao (chia đều chi phí của thao tác đắt đỏ cho chuỗi các thao tác $\rightarrow$ ví dụ: DSU, SegTree).
    \item[Lexicographically largest] Lớn nhất theo thứ tự từ điển.
    \item[Pigeonhole Principle] Nguyên lý chuồng bồ câu (Dirichlet - bẫy toán tư duy đếm hoặc tìm chu trình lặp).
    \item[Heuristic] Thuật toán xấp xỉ / Phỏng đoán (không đảm bảo tối ưu 100\% nhưng chạy nhanh, dùng trong bài toán tối ưu NP-hard).
    \item[Alternately] Một cách luân phiên (hết người/thao tác này rồi bắt buộc tới người/thao tác kia).
\end{description}

% \columnbreak

\vspace{2.25cm}
\hrule
\vspace{0.25cm}

\begin{center}
\textbf{TRÌNH - HIEUTHUHAI}
\end{center}

\begin{multicols*}{2}
\raggedright
\small

\textbf{Tiếng Việt} \\
\vspace{0.3em}
Ey ey yeah yeah \\
Ối zồi ôi ối zồi ôi \\
Trình là gì mà là trình ai chấm \\
Anh chỉ biết làm ba mẹ tự hào \\
Xây căn nhà thật to ở một mình hai tấm \\
Ối zồi ôi ối zồi ôi \\
Cứ lên mạng phán xét tưởng là mình oai lắm \\
Nhìn vào sự nghiệp anh thèm chảy nước miếng \\
Giống mấy thằng biến thái nó đang rình ai tắm \\
Mua bao nhiêu căn nhà (đếm được đếm được) \\
Năm bao nhiêu show (đếm được đếm được) \\
Bao nhiêu bài hit (đếm được đếm được) \\
Bao nhiêu trong bank (đếm được đếm được) \\
Mua bao nhiêu căn nhà (đếm được đếm được) \\
Năm bao nhiêu show (đếm được đếm được) \\
Bao nhiêu là cúp (đếm được đếm được) \\
Có bao nhiêu fan (không đếm được không đếm được)

\columnbreak

\textbf{English} \\
\vspace{0.3em}
Ey ey yeah yeah \\
Oh my god oh my god \\
What is skill, and who has the right to judge? \\
I only know how to make my parents proud \\
Building a huge 2-story house to live by myself \\
Oh my god oh my god \\
Keep judging online, thinking you're so cool \\
Staring at my career, drooling with envy \\
Just like perverts peeping at someone taking a bath \\
How many houses bought? (countable, countable) \\
How many shows a year? (countable, countable) \\
How many hit songs? (countable, countable) \\
How much in the bank? (countable, countable) \\
How many houses bought? (countable, countable) \\
How many shows a year? (countable, countable) \\
How many trophies? (countable, countable) \\
How many fans do I have? (uncountable, uncountable!)

\end{multicols*}
